\documentclass[]{article}  

% Default font is the helvetica postscript font
\usepackage{helvet}
\usepackage{hyperref}
\usepackage{amsmath}
\usepackage{amsthm}
\usepackage{amssymb}
\usepackage[a4paper,margin=25mm]{geometry}
\usepackage{fancyhdr}
\usepackage{lipsum}
%\usepackage[ruled,vlined]{algorithm2e}
\usepackage[utf8]{inputenc}
\usepackage[english]{babel}
\usepackage{graphicx}


\newtheorem{theorem}{Theorem}
\newtheorem{lemma}[theorem]{Lemma}
\newtheorem{defn}[theorem]{Definition}
\newtheorem{remark}[theorem]{Remark}
\newtheorem{corollary}[theorem]{Corollary}
\newtheorem{note}[theorem]{Note}

%\renewcommand{\baselinestretch}{1.15}
\newcommand{\s}{\mathbf{s}}
\newcommand{\ba}{\mathbf{a}}
\newcommand{\bw}{\mathbf{w}}
\newcommand{\bx}{\mathbf{x}}
\newcommand{\bv}{\mathbf{v}}
\newcommand{\thetab}{\boldsymbol{{\theta}}}
\newcommand{\h}{\mathbf{h}}

\newcommand{\cA}{\mathcal{A}}
\newcommand{\cW}{\mathcal{W}}
\newcommand{\cH}{\mathcal{H}}

\newcommand{\R}{\mathbb{R}}
\newcommand{\E}{\mathbb{E}}
\newcommand{\Prob}{\mathbb{P}}
\newcommand{\KL}{\mathrm{D}_{KL}}



\title{On the Maxima of Sub-Gaussian Random Variables}
\author{Behrad Moniri\\\texttt{bemoniri@live.com}}
\date{}

\usepackage{sectsty} \sectionfont{\fontsize{15}{15}\selectfont}
\begin{document}
	\maketitle
	
	
	\begin{abstract}
		In this note, we will first prove a	bound on the expected maxima of a sequence of weighted sub-gaussian random variables. Next, an upper bound for the expected value of the maximum of a finite number of sub-gaussian random variables is proved and also a high probability version of this statement is presented.
	\end{abstract}
	
	
	The following theorem on the expected maxima of an infinite sequence of weighted sub-gaussian random variables is stated in Excersise (2.5.10) of \cite{vershynin}.
	\begin{theorem}
		Let $X_1, X_2, \dots$ be a sequence of independent $\sigma$-sub Gaussian random variables, then
		\begin{equation}
		\label{bound}
			\E \Bigg[\max_i \frac{|X_i|}{\sqrt{1+\log(i)}}\Bigg] \leq 2\sigma + \frac{\pi^2}{3}\sqrt{2\pi}.
		\end{equation}
	\end{theorem}
	
	\begin{proof}
		The expected value can be written in terms of an integral of tail probabilities:
		
		\begin{equation*}
			\E\Bigg[\frac{|X_i|}{\sqrt{1+\log(i)}}\Bigg] = \int_{0}^{\infty} \mathbb{P}\Big(\max_i \frac{|X_i|}{\sqrt{1+\log(i)}} \geq t\Big) dt
		\end{equation*}
		
		If $X$ is $\sigma$-sub gaussian, we have $\Prob[|X|\geq t] \leq 2e^{-\frac{t^2}{2}}$.
		Let $a = 2\sigma$, we can divide the integral into two parts and write the following chain of inequalities
		
		\begin{align*}
		\E\Bigg[\frac{|X_i|}{\sqrt{1+\log(i)}}\Bigg] &\leq \int_{0}^{2\sigma}
		dt	
		+
		\int_{2\sigma}^{\infty}		
		 \mathbb{P}\Bigg(\max_i \frac{|X_i|}{\sqrt{1+\log(i)}} \geq t\Bigg)\; dt\\		 
		&\leq 2\sigma + \int_{2\sigma}^{\infty} \sum_{i = 1}^{\infty} \Prob\Bigg(\frac{|X_i|}{\sqrt{1+\log(i)}}\Bigg)\; dt && \text{(union bound)}\\
		&\leq 2\sigma + \sum_{i = 1}^{\infty} \int_{2\sigma}^{\infty} 
		2\exp\Big(\frac{-t^2}{2\sigma^2} \big(1+\log(i)\big)\Big)
		\; dt && \text{(sub-gaussian tails)}\\		
		&\leq 2\sigma + 2\sum_{i = 1}^{\infty} \int_{2}^{\infty} 
		e^{-u^2/2}.\;i^{-u^2/2}
		\; du \\
		&\leq 2\sigma + 2\sum_{i = 1}^{\infty} \int_{2}^{\infty} 
		e^{-u^2/2}.\;i^{-2}
		\; du \\
		&\leq 2\sigma + 2\sum_{i = 1}^{\infty}i^{-2} \int_{2}^{\infty} 
		e^{-u^2/2}
		\; du \\
		&\leq 2\sigma + \frac{\pi^2}{3}\sqrt{2\pi}.
		\end{align*}		
	\end{proof}

\noindent
We will now prove a similar bound, for the maximum of finite number of sub-gaussian random variables.

\begin{theorem}
	\label{bound2}
	Let $X_1, \dots, X_n$ be independent $\sigma$-sub gaussian random variables. We have
	\begin{equation}
		\E \big[\max_{i\in\{1, \dots, n\}} X_i\big] \leq \sigma\sqrt{2\log n}.
	\end{equation}
\end{theorem}

\begin{proof}
	Let $Y =  \max_{i\in\{1, \dots, n\}} X_i$.
	\begin{align*}
		e^{\lambda \E[\max X_i]} &\leq \E\Big[e^{\lambda \max X_i}\Big] \\
		&= \E\Big[\max e^{\lambda X_i}\Big]\\
		&\leq \E\Bigg[\sum_{i = 1}^{n} e^{\lambda X_i} \Bigg]\\
		&\leq n e^{\frac{\lambda^2\sigma^2}{2}}.
	\end{align*}
	Where the first inequality is a consequence of Jensen's inequality.
	Hence, 
	\begin{equation*}
		\E[\max X_i] \leq \frac{\log(n)}{\lambda} + \frac{\lambda \sigma^2}{2}.
	\end{equation*}
	Optimizing the RHS, yields
	$\lambda^* = \sqrt{\frac{2\log n}{\sigma^2}}$. Thus
	\begin{equation*}
		\E[\max X_i] \leq \sigma \sqrt{2\log(n)}
	\end{equation*}
	
	
\end{proof}


\begin{note}
	Let $X_1, \dots, X_n \sim \mathcal{N}(0, \sigma^2)$ be i.i.d. random variables. In \cite{gautam}, the following upper and lower bounds are proved:
	\begin{equation}
		\frac{1}{\sqrt{\pi\log 2}} \sigma \sqrt{\log n} \leq \E \big[\max_{i\in\{1, \dots, n\}} X_i\big] \leq  \sigma \sqrt{2\log n}.
	\end{equation}		
	Hence, the bound \eqref{bound2} is sharp.
\end{note}


\noindent
We will now show that the maximum, is less than $\sqrt{2\sigma^2\log n}$ with high probability.

\begin{theorem}
	Let $X_1, \dots, X_n$ be independent $\sigma$-sub gaussian random variables:
	\begin{equation}
		\mathbb{P}\Big(\max X_i -\sqrt{2\sigma^2\log n} \geq t\Big) \leq \exp\bigg(\frac{-t^2}{2\sigma^2}\bigg).
	\end{equation}
\end{theorem}


\begin{proof}
	\begin{align*}
		\mathbb{P}\Big(\max X_i -\sqrt{2\sigma^2\log n} \geq t\Big) &= \mathbb{P}\Big( \exists i\in[n]:\;\; X_i \geq t + \sqrt{2\sigma^2\log n} \Big)\\		
		&\leq n\; \mathbb{P}\Big( X_i \geq t + \sqrt{2\sigma^2\log n} \Big) && \text{(union bound)}\\
		&\leq n\; \exp\bigg(\frac{-\big(t + \sqrt{2\sigma^2\log n}\big)^2}{2\sigma^2}\bigg) && \text{(sub-gaussian tail)}\\
		&= \exp\bigg(\frac{-t^2}{2\sigma^2}\bigg)\exp\bigg(\frac{-2t \sqrt{2\sigma^2\log n}}{2\sigma^2}\bigg)\leq \exp\bigg(\frac{-t^2}{2\sigma^2}\bigg).
	\end{align*}
	
\end{proof}

\bibliography{bib}
\bibliographystyle{alpha}
		
\end{document}